\scanpage{62}
\deleted{\unsure{un fonc}} et que l'hom naturel $\underline\Gamma(W_T/T)\unsure{-|\to}\underline\Gamma(W_{T'}/T')$ pour $T'\to T$ \deleted{\unsure{de fa|défini}} soit \uline{local} sur les fibres en faisceaux. On obtient ainsi un foncteur \inserted{$\unsure{\underline{\mathbb W}|\underline{\underline{\mathbb W}}}$} de nature \unsure{locale|covt|covariant} de $S$-préschémas en $S$-préschémas, et un isom fonctoriel $|\underline{\mathbb W}(T)|\simeq|T|$ (espaces ss-jacents). \deleted{\unsure{on peut évidemment|on voit évidemment|vraisemblablement} d'ailleurs $\Lambda$, puisque \unsure{l'on aurait} $\mathrm{Hom}_S(T,\Lambda)=$} \inserted{Nous donnerons des} conditions suffisantes plus \unsure{loin|bas} pour que $\Lambda$ ait les \unsure{propr. voulues|voisinage}.

\uline{Exemple} III Soit $S'=\operatorname{Spec}(A)$, où $A$ est un faisceau loc. libre d'algèbres comm. sur $\underline O_S$. Supposons de plus que $S'\unsure{-|\to}S$ soit un homéomorphisme universel. Posons $\Lambda=\deleted{\unsure{\operatorname{Spec}}}\mathbf V(\check A)$. On est alors simultanément dans le cas de l'exemple I et de l'exemple II.
\margin{$S'\to S$}

\uline{Exemple} III' Comme II, mais on suppose donnée une augmentation $W\longrightarrow S[t]$ d'anneaux, telle que pour tout $T/S$, l'augmentation correspondante $\underline O_{\underline{\mathbb W}(T)}\longrightarrow\underline O_T$ soit surjective, de noyau nilpotent.
% slanted
\margin{\unsure{blague sur} exemple II : foncteurs de Greenberg}

Soit $\unsure{\mathfrak X|\mathcal X|\underline X}$ un préschéma sur $\underline{\mathbb W}(S)$. Considérons le foncteur contravariant en le $S$-préschéma $T$
\[
\mathrm{Hom}_{\mathbb W(S)}(\mathbb W(T),\mathfrak X)\tag*{(*)}
\]
% archive sheet number: 61
